\documentclass[10pt,a4paper]{amsart}
\usepackage[margin=19mm]{geometry}
\usepackage{amsmath,amssymb,mathtools}
\usepackage[scr=rsfs,cal=euler]{mathalfa}
\usepackage{xfrac}
\usepackage[unicode,hidelinks,bookmarks=false]{hyperref}
\newcommand{\ie}{i.e.\ }
\newcommand{\cf}{cf.\ }
\newcommand{\resp}{respectively\ }
\newcommand{\Subsection}{\subsection}
\providecommand{\nobreakdash}{\nobreak\hbox{-}}
\setlength{\emergencystretch}{2em}
\multlinegap0pt
\usepackage{bm}
\usepackage{bbm}
\usepackage{dsfont}
\usepackage{xspace}
\usepackage{cancel}
\usepackage{mathtools}
\usepackage{amsthm}
\makeatletter
\@ifundefined{thm}{\newtheorem{thm}{Theorem}}{}
\makeatother
\title[Fuchsian differential equations of order 3,...,6 with three ]{Fuchsian differential equations of order 3,...,6 with three singular points and an accessory parameter}
\author{Yoshishige Haraoka, Hiroyuki Ochiai, Takeshi Sasaki, Masaaki Yoshida}
\date{Source: arXiv:2307.01358v3 (CC BY 4.0)}
\begin{document}
\maketitle

\noindent\emph{Source and licence.} Fuchsian differential equations of order 3,...,6 with three singular points and an accessory parameter. Version arXiv:2307.01358v3 (CC BY 4.0), distributed under the Creative Commons Attribution 4.0 International licence, as recorded in the arXiv licence metadata for this exact version and shown on the abstract page https://arxiv.org/abs/2307.01358v3. Full rights notice: licenses/arxiv-2307.01358-CC-BY-4.0.txt.\par\smallskip

\noindent\textit{Editorial scope.} Counted unit: the Thm whose source label is \texttt{red\_atoap} in the article (source0.tex lines 1762--1765, source label \texttt{red\_atoap}), delivered together with the article's own complete original proof of it (source0.tex lines 1766--1794, one \texttt{proof} environment of 29 lines). No mathematics is written, repaired or re-derived: statement and proof are the article's own text line for line. Required context retained uncounted: none. Every definition, display and label of those lines is retained unchanged. Omitted from this package and not redistributed: 1--1761, 1795--3329. Nothing inside a retained excerpt is omitted or altered: every retained body is delivered exactly as the article's lines are written, so no omission receipt of any kind accompanies this package. Citations: the retained excerpts carry no citation command at all, so no bibliography entry of the article is transcribed and no reference is redistributed. Reference adaptation: none was needed and none was made; every reference inside the delivered unit resolves inside it, so no unresolved reference or citation is produced. Printed slips kept literal: no printed word, symbol, hypothesis or number of the counted statement or of its proof was altered. Layout and class: the article's own class is \texttt{article} with packages including amsmath, bm, amssymb, graphicx, setspace, enumitem, lscape, amsthm, minitoc, color; the wrapper is the lane's pinned \texttt{amsart} editorial wrapper at 10pt on A4, which restates the theorem-like environments the retained text needs and adds the article's own macro definitions that the retained text needs (none), with extra wrapper packages (bm, bbm, dsfont, xspace, cancel, mathtools). The delivered numbering is the unit's own. Assets: no retained span contains a figure environment, a graphics-inclusion command, a PDF-page inclusion, a table or a TikZ picture; no image file is copied into this package, the package declares no asset dependency and no image file is redistributed with it. Licence: version arXiv:2307.01358v3 of this article is distributed under the Creative Commons Attribution 4.0 International licence, as recorded in the arXiv licence metadata for this exact version and shown on the abstract page https://arxiv.org/abs/2307.01358v3. A full rights notice is delivered at licenses/arxiv-2307.01358-CC-BY-4.0.txt. Characters: every retained excerpt is pure ASCII, so the delivered engine, XeLaTeX with its default Latin Modern text font, reports no missing character.
\begin{thm}\label{red_atoap}
  Assume $H$ and $H'$ are connected by the shift relation $H'P=QH$.
  If $H$ is reducible, so is $H'$. If $H'$ is reducible, so is $H$.
\end{thm}

\begin{proof}Assume $H$ is reducible:
  \[H=F_1\circ F_2,\quad n_j={\rm order}(F_j),\quad j=1,\, 2,\]
and $F_2$ is irreducible. Then, considering 
the dimension of $P({\rm Sol}(F_2))$, we have three cases:

\smallskip
(1)\quad $\dim P({\rm Sol}(F_2))=n_2$,

\smallskip
(2) \quad $0<\dim P({\rm Sol}(F_2))<n_2$,

\smallskip
(3)\quad $P({\rm Sol}(F_2))=0$.
\smallskip

In the first case, 
$H'$ has an $n_2$-dimensional solution space $P({\rm Sol}(F_2))$, and, therefore, it is divisible by an irreducible operator of order $n_2$. Thus $H'$ is reducible.

The second case does not occur because
the kernel of $P$ is a nontrivial invariant subspace of ${\rm Sol}(F_2)$ and this contradicts to the irreducibility of $F_2$.

Assume the third case; we write $P$ as $P=P_1\circ F_2$
and divide both sides of $H'P=QH$ by $F_2$. Then, we have
  \[H'\circ P_1=Q\circ F_1.\]
Since ${\rm order} (P)<n=n_1+n_2$, we see that
${\rm order} (P_1) < n_1$ and that $P_1({\rm Sol}(F_1))\not=0$.
Thus ${\rm Sol}(H')$ admits a non-trivial invariant subspace,
which implies that $H'$ is reducible.
The latter statement is obtained by taking adjoint. \end{proof}

\end{document}
